\section{Stieltjes convolution identities and repeated primitives}
\label{per:sec:convolutions}

\subsection{All Stieltjes functions as convolution polynomials}

The regularized germ
\begin{equation}\label{per:eq:B}
 B(u)=-uZ_{1+u}=\Gamma(1-u)W_{1+u},\qquad |u|<1,
\end{equation}
satisfies
\begin{equation}\label{per:eq:Bseries}
 B(u)=\Delta+
       \sum_{r\ge1}\frac{(-1)^r}{(r-1)!}T_{r-1}u^r.
\end{equation}
The following formula is a statement in the commutative convolution
algebra with identity $\Delta$; exponential and Bell-polynomial
expressions can equivalently be read as formal power series.

\begin{theorem}[Stieltjes convolution generator]
\label{per:thm:Bell}
For $|u|<1$,
\begin{equation}\label{per:eq:exp}
 B(u)=\exp_*\left(uP_0+
                    \sum_{r\ge2}\frac{\zeta(r)}{r}u^r\Delta\right).
\end{equation}
Here the series defines the analytic germ after pairing with any smooth
test function. In particular, if $\mathcal B_j^*$ denotes the complete
exponential Bell polynomial with every multiplication interpreted as
convolution and every scalar constant as that scalar times $\Delta$,
then
\begin{equation}\label{per:eq:Bell}
 T_n=\frac{(-1)^{n+1}}{n+1}
 \mathcal B_{n+1}^*
 \bigl(P_0,\zeta(2)\Delta,2!\zeta(3)\Delta,
                         \ldots,n!\zeta(n+1)\Delta\bigr).
\end{equation}
The two-parameter form is
\begin{equation}\label{per:eq:Bconv}
 B(u)*B(v)=C(u,v)B(u+v),
 \qquad C(u,v)=\frac{\Gamma(1-u)\Gamma(1-v)}{\Gamma(1-u-v)},
\end{equation}
as an analytic germ at $(u,v)=(0,0)$. Its scalar correction has
\begin{equation}\label{per:eq:C}
 \log C(u,v)=\sum_{r\ge2}\frac{\zeta(r)}{r}
                   \bigl(u^r+v^r-(u+v)^r\bigr).
\end{equation}
Thus every $T_m*T_n$ is a finite explicit linear combination of
$\Delta,T_0,\ldots,T_{m+n+1}$, with coefficients polynomial in positive
integer zeta values.
\end{theorem}

\begin{proof}
At a nonzero Fourier mode, $\widehat P_0(n)=\gamma+L_n$. The Taylor
expansion
\[
 \log\Gamma(1-u)=\gamma u+
                     \sum_{r\ge2}\frac{\zeta(r)}r u^r
\]
therefore proves \eqref{per:eq:exp} mode by mode. The same polynomial
growth argument as in Theorem~\ref{per:thm:family} proves normal
convergence in every smaller disk after test-function pairing. Comparing
the complete Bell-polynomial generating function with
\eqref{per:eq:Bseries} gives \eqref{per:eq:Bell}; the factorial is
$n!/(n+1)!=1/(n+1)$. Equation \eqref{per:eq:Bconv} follows from
\eqref{per:eq:group}. The logarithm of its Gamma quotient is
\eqref{per:eq:C}; its linear Euler-constant terms cancel. Coefficient
extraction at $u^{m+1}v^{n+1}$ gives the last assertion.
\end{proof}

The first identities are
\begin{align}
 T_0&=-P_0,\notag\\
 T_1&=\frac12\bigl(P_0^{*2}+\zeta(2)\Delta\bigr),
 \label{per:eq:T1}\\
 T_2&=-\frac13\bigl(P_0^{*3}+3\zeta(2)P_0+2\zeta(3)\Delta\bigr),
 \label{per:eq:T2}\\
 T_3&=\frac14\bigl(P_0^{*4}+6\zeta(2)P_0^{*2}
          +8\zeta(3)P_0+(3\zeta(2)^2+6\zeta(4))\Delta\bigr).
 \label{per:eq:T3}
\end{align}
Here $P_0^{*r}$ is an $r$-fold convolution, not a pointwise power.
Eliminating these powers gives, for example,
\begin{align}
 T_0*T_0&=2T_1-\zeta(2)\Delta,\label{per:eq:00}\\
 T_0*T_1&=\frac32T_2-\zeta(2)T_0+\zeta(3)\Delta,
 \label{per:eq:01}\\
 T_0*T_2&=\frac43T_3-2\zeta(2)T_1+2\zeta(3)T_0
                                      -2\zeta(4)\Delta,
 \label{per:eq:02}\\
 T_1*T_1&=T_3-2\zeta(2)T_1+2\zeta(3)T_0
                                      -\frac14\zeta(4)\Delta.
 \label{per:eq:11}
\end{align}
The last coefficient uses $\zeta(2)^2=\tfrac52\zeta(4)$.

\subsection{The reflected Stieltjes generator}

Define the periodic Hilbert transform on zero-mean distributions by
\begin{equation}\label{per:eq:Hilbert}
 \widehat{\mathcal H U}(n)=-i\operatorname{sgn}(n)\widehat U(n)
 \quad(n\ne0),\qquad \widehat{\mathcal H U}(0)=0.
\end{equation}
Its multiplier is bounded, so it is a continuous operator on periodic
distributions and commutes with the order coefficient operations above.
Reflection in the following theorem acts on the periodic variable only;
it does not conjugate the complex parameters.

\begin{theorem}[All-order reflected convolution]
\label{per:thm:reflectedgenerator}
As an analytic germ at $(u,v)=(0,0)$,
\begin{equation}\label{per:eq:Brefconv}
 B(u)*\check{B(v)}
 =C(u,v)\bigl(\cos(\pi v)\,\mathrm{Id}
                    +\sin(\pi v)\,\mathcal H\bigr)B(u+v),
\end{equation}
with the same Gamma quotient $C(u,v)$ as in \eqref{per:eq:Bconv}.
In particular,
\begin{align}
 T_0*\check T_1={}&\tfrac12T_2+\check T_2
                +\zeta(2)(T_0+\check T_0)+\zeta(3)\Delta,
 \label{per:eq:T0refT1}\\
 T_1*\check T_1={}&\tfrac12(T_3+\check T_3)
                +2\zeta(2)(T_1+\check T_1)
                +\zeta(3)(T_0+\check T_0)
                +\tfrac72\zeta(4)\Delta.
 \label{per:eq:T1refT1}
\end{align}
More generally, every $T_m*\check T_n$ reduces to a finite linear
combination of $\Delta$, $T_j$ and $\check T_j$ with
$0\le j\le m+n+1$ and coefficients polynomial in positive integer
zeta values over $\mathbb Q$.
\end{theorem}

\begin{proof}
Since $L_{-n}=L_n-i\pi\operatorname{sgn}n$, the nonzero Fourier
coefficient of the left side of \eqref{per:eq:Brefconv} is
\[
 \Gamma(1-u)\Gamma(1-v)
       e^{(u+v)L_n}e^{-i\pi v\operatorname{sgn}n}.
\]
The last factor equals
$\cos(\pi v)-i\operatorname{sgn}(n)\sin(\pi v)$, precisely the
multiplier of the operator on the right. The normal-convergence argument
for the normalized family proves the identity as a distribution-valued
analytic germ and justifies every fixed mixed coefficient extraction.
This also checks the sign of the Hilbert-transform term directly.

For explicit coefficients put $b=\gamma+\log(2\pi|n|)$ and
$a=\tfrac{\pi i}{2}\operatorname{sgn}n$, so $a^2=-\pi^2/4$.
Equations \eqref{per:eq:T1}--\eqref{per:eq:T3} give polynomials
$t_j(y)=\widehat T_j(n)$ in $y=b+a$, while reflection replaces $a$
by $-a$. Direct polynomial multiplication gives
\begin{align*}
 t_0(b+a)t_1(b-a)
 ={}&\tfrac12t_2(b+a)+t_2(b-a)
    +\zeta(2)\bigl(t_0(b+a)+t_0(b-a)\bigr)+\zeta(3),\\
 t_1(b+a)t_1(b-a)
 ={}&\tfrac12\bigl(t_3(b+a)+t_3(b-a)\bigr)
    +2\zeta(2)\bigl(t_1(b+a)+t_1(b-a)\bigr)\\
 &+\zeta(3)\bigl(t_0(b+a)+t_0(b-a)\bigr)+\tfrac72\zeta(4).
\end{align*}
Here $\pi^2=6\zeta(2)$ and $\pi^4=90\zeta(4)$ have been used.
These prove the two displayed identities at every nonzero Fourier mode;
all zero modes vanish.

Finally $t_j$ has degree $j+1$ and nonzero rational leading coefficient
$(-1)^{j+1}/(j+1)$. Put $M=m+n+2$. After reducing $a^2=-\pi^2/4$,
the product $t_m(b+a)t_n(b-a)$ has the form
$A(b)+aE(b)$ with $\deg A\le M$ and $\deg E\le M-1$.
The symmetric polynomials $t_j(b+a)+t_j(b-a)$, $0\le j<M$,
have successive degrees $1,\ldots,M$ in $b$, while their antisymmetric
differences divided by $a$ have successive degrees $0,\ldots,M-1$.
Together with the constant they form a triangular basis of this space.
Elimination divides only by nonzero rational leading coefficients.
Its coefficients belong to
$\mathbb Q[\pi^2,\zeta(2),\ldots,\zeta(M)]$, which is the asserted
zeta algebra because $\pi^2=6\zeta(2)$.
\end{proof}

The frequency-sign contribution can equivalently be encoded by
\[
 P_0-\check P_0=-\pi\mathcal H\Delta,\qquad
 (P_0-\check P_0)^{*2}=-\pi^2\Delta.
\]
Thus reflection adds a two-dimensional sign algebra to the unreflected
Bell calculus, with its constant term fixed by the same endpoint mass.

\subsection{Polygamma correlations with all contact terms}

\begin{theorem}[Convolution and reflected convolution]
\label{per:thm:polycorr}
Let $k,l\ge0$ and $r=k+l$. With the finite-part normalization
\eqref{per:eq:Pdef},
\begin{align}
 P_k*P_l=D^r\bigl\{&2T_1+(H_k+H_l)T_0
                         +(H_kH_l-\zeta(2))\Delta\bigr\},
 \label{per:eq:PP}\\
 P_k*\check P_l=(-1)^lD^r\bigl\{&T_1+\check T_1
                   +H_lT_0+H_k\check T_0
                   +(H_kH_l+2\zeta(2))\Delta\bigr\}.
 \label{per:eq:PrefP}
\end{align}
For the derivative-compatible family the harmonic coefficients cancel:
\begin{align}
 Q_k*Q_l&=D^r\{2T_1-\zeta(2)\Delta\},
 \label{per:eq:QQ}\\
 Q_k*\check Q_l&=(-1)^lD^r
              \{T_1+\check T_1+2\zeta(2)\Delta\}.
 \label{per:eq:QrefQ}
\end{align}
\end{theorem}

\begin{proof}
Put $y_n=\gamma+L_n$. The first formula follows from
\[
 \widehat T_0(n)=-y_n,\qquad
 2\widehat T_1(n)=y_n^2+\zeta(2)
\]
and $(y_n-H_k)(y_n-H_l)=y_n^2-(H_k+H_l)y_n+H_kH_l$.
For reflection set $b_n=\gamma+\log(2\pi|n|)$ and
$a_n=\frac\pi2\operatorname{sgn}n$, so $y_n=b_n+ia_n$ and
$y_{-n}=b_n-ia_n$. Then
\[
 (y_n-H_k)(y_{-n}-H_l)
 =b_n^2+a_n^2-(H_k+H_l)b_n+H_kH_l
                                      +ia_n(H_k-H_l).
\]
Also
\[
 \widehat T_1(n)+\widehat T_1(-n)=b_n^2-a_n^2+\zeta(2),
\]
and
$H_l\widehat T_0(n)+H_k\widehat T_0(-n)
=-(H_k+H_l)b_n+ia_n(H_k-H_l)$.
The missing constant is $2a_n^2-\zeta(2)=2\zeta(2)$.
Finally $(2\pi i(-n))^l=(-1)^l(2\pi in)^l$ accounts for the
prefactor in \eqref{per:eq:PrefP}. The last two identities follow
from $Q_k=D^kP_0$ and $\check{D^lP_0}=(-1)^lD^l\check P_0$.
\end{proof}

At the first level,
\begin{equation}\label{per:eq:reflected0}
 P_0*\check P_0=T_1+\check T_1+
                         \frac{\pi^2}{3}(\delta_0-1).
\end{equation}
For $0<a<1$ its ordinary representative is
$\gamma_1(a)+\gamma_1(1-a)-\pi^2/3$. Equation
\eqref{per:eq:reflected0} supplies the missing endpoint mass in the
formal twice-differentiated log-Gamma correlation. In particular, if
$C_g=g*\check g$, then
\begin{equation}\label{per:eq:Cg}
 -D^2C_g=T_1+\check T_1+\frac{\pi^2}{3}(\delta_0-1).
\end{equation}
The negative sign follows from differentiating the reflected factor.

For $r\ge1$, the representative of \eqref{per:eq:PP} on $(0,1)$ can
also be written with one positive-order Hurwitz jet:
\begin{equation}\label{per:eq:PPordinary}
 (P_k*P_l)(a)=(-1)^{r+1}r!
 \left\{2\zeta'(r+1,a)+(2H_r-H_k-H_l)\zeta(r+1,a)\right\}.
\end{equation}
Indeed, differentiating the Hurwitz shift identity at $s=1$ gives
$\gamma_1^{(r)}(a)=(-1)^{r+1}r![\zeta'(r+1,a)+H_r\zeta(r+1,a)]$.
For $r=0$, the separate constant $\zeta(2)$ from $-\zeta(2)\Delta$
must be retained.

Translations cause no further ambiguity. If
$\tau_aU(x)=U(x-a)$ for real $a$ modulo $1$, then
$(\tau_aU)*(\tau_bV)=\tau_{a+b}(U*V)$ and
$\tau_aU*\check{\tau_bV}=\tau_{a-b}(U*\check V)$.
All contact terms above are therefore translated to the sum or difference
of the two shifts. General complex translations are not operators on all
periodic distributions, because their Fourier multipliers grow
exponentially in one direction; they require a separately specified
analytic or one-sided boundary-value setting.

\subsection{Convergent integral identities without distribution notation}

The distribution formulas have ordinary, absolutely convergent integral
versions at every nonintegral shift. The following construction makes
the subtraction completely explicit. For $m,n\ge0$ and $0<a<1$, put
\begin{align}
 \mathcal I_{m,n}(a)={}&
 \int_0^a\left\{\gamma_m(x)\gamma_n(a-x)
       -\gamma_n(a)\frac{\log^m x}{x}
       -\gamma_m(a)\frac{\log^n(a-x)}{a-x}\right\}dx\notag\\
 &+\int_a^1\gamma_m(x)\gamma_n(1+a-x)\,dx\notag\\
 &+\frac{\gamma_n(a)}{m+1}\log^{m+1}a
       +\frac{\gamma_m(a)}{n+1}\log^{n+1}a.
 \label{per:eq:Imn}
\end{align}

\begin{proposition}[Ordinary integral realization]
\label{per:prop:Imn}
Both integrals in \eqref{per:eq:Imn} converge absolutely, and
\[
 \mathcal I_{m,n}(a)=(T_m*T_n)(a),\qquad 0<a<1.
\]
Consequently,
\begin{align}
 \mathcal I_{0,0}(a)&=2\gamma_1(a)+\zeta(2),
 \label{per:eq:I00}\\
 \mathcal I_{0,1}(a)&=\frac32\gamma_2(a)-\zeta(2)\gamma_0(a)-\zeta(3),
 \label{per:eq:I01}\\
 \mathcal I_{0,2}(a)&=\frac43\gamma_3(a)-2\zeta(2)\gamma_1(a)
                         +2\zeta(3)\gamma_0(a)+2\zeta(4),
 \label{per:eq:I02}\\
 \mathcal I_{1,1}(a)&=\gamma_3(a)-2\zeta(2)\gamma_1(a)
                         +2\zeta(3)\gamma_0(a)+\frac14\zeta(4).
 \label{per:eq:I11}
\end{align}
\end{proposition}

\begin{proof}
Near $x=0$, the shift formula for $\gamma_m$ and smoothness of
$\gamma_n$ near $a$ show that the first subtracted integrand is
$O(1+|\log x|^m)$. The same argument applies near $x=a$. The second
integral has no singular endpoints for fixed $a\in(0,1)$.

For a direct identification, first restrict $a$ to a compact subinterval
of $(0,1)$. The singular supports of $T_m(x)$ and $T_n(a-x)$ are
separated in the integration variable. A smooth partition of unity
around $x=0$, $x=a$ and their complement therefore defines their
product: in each singular neighborhood only one factor is a
distribution and the other is smooth. Applying the two cutoff formulas
\eqref{per:eq:Tcutoff}, and then combining the complementary integrals,
gives
\begin{align*}
 \lim_{\epsilon\downarrow0}\biggl\{&
 \int_\epsilon^{a-\epsilon}\gamma_m(x)\gamma_n(a-x)\,dx
 +\int_a^1\gamma_m(x)\gamma_n(1+a-x)\,dx\\
 &+\frac{\gamma_n(a)}{m+1}\log^{m+1}\epsilon
  +\frac{\gamma_m(a)}{n+1}\log^{n+1}\epsilon\biggr\}.
\end{align*}
The use of cutoff singular coefficients at $a$ introduces errors at most
$O(\epsilon|\log\epsilon|^M)$ for a fixed $M$, hence no additional
constant. Integrating the explicitly subtracted logarithmic terms
evaluates this limit as \eqref{per:eq:Imn}. This agrees with convolution
defined by addition on the torus. Equations \eqref{per:eq:I00}--
\eqref{per:eq:I11} follow by restricting \eqref{per:eq:00}--
\eqref{per:eq:11} to $(0,1)$, where $\Delta$ is the constant $-1$.
\end{proof}

In terms of the ordinary digamma function, the first identity reads
\begin{align}
 2\gamma_1(a)+\zeta(2)={}&
 \int_0^a\left\{\psi(x)\psi(a-x)
       +\psi(a)\left(\frac1x+\frac1{a-x}\right)\right\}dx\notag\\
 &+\int_a^1\psi(x)\psi(1+a-x)\,dx-2\psi(a)\log a.
 \label{per:eq:digammaintegral}
\end{align}
No divergent integral occurs on the right: the braces are to be combined
before integration. For instance,
\[
 \mathcal I_{0,0}(1/2)
 =2\gamma_1-4\gamma\log2-2\log^22+\zeta(2),
\]
because $\gamma_1(1/2)=\gamma_1-2\gamma\log2-\log^22$.
The factor of the unshifted $\gamma_1$ in the latter identity is $1$,
as follows directly from $\zeta(s,1/2)=(2^s-1)\zeta(s)$.

\subsection{Every repeated Stieltjes primitive in one finite formula}

For a zero-mean periodic distribution define $D^{-q}$ by its Fourier
multiplier $(2\pi i n)^{-q}$ at $n\ne0$ and zero at $n=0$. This is
the unique repeated primitive having zero mean at each stage; it equals
convolution with $W_{1-q}$. Define the rational coefficients
\begin{equation}\label{per:eq:h}
 \prod_{r=1}^{q-1}\left(1-\frac wr\right)^{-1}
       =\sum_{j\ge0}h_j^{(q-1)}w^j,
 \qquad h_0^{(0)}=1,\quad h_j^{(0)}=0\ (j>0).
\end{equation}
They are complete homogeneous symmetric functions in
$1,1/2,\ldots,1/(q-1)$; equivalently their generating function is
\[
 \exp\left(\sum_{r\ge1}\frac{H_{q-1}^{(r)}}r w^r\right),
 \qquad H_N^{(r)}=\sum_{j=1}^N\frac1{j^r}.
\]
For example $h_1^{(N)}=H_N$ and
$h_2^{(N)}=(H_N^2+H_N^{(2)})/2$.

\begin{theorem}[All repeated primitives and their constants]
\label{per:thm:primitive}
For $q\ge1,n\ge0$, the zero-mean $q$-fold primitive of $T_n$ is
\begin{equation}\label{per:eq:primitive}
 D^{-q}T_n(x)=\frac{(-1)^{n+1}n!}{(q-1)!}
 \sum_{j=0}^{n+1}\frac{h_{n+1-j}^{(q-1)}}{j!}
                              \zeta^{[j]}(1-q,x),\quad 0<x<1.
\end{equation}
Here $\zeta^{[j]}$ means the $j$th derivative in the first argument.
The right side is periodized as a locally integrable function. For
$q\ge2$ it is continuous, with the regularity implied by its endpoint
terms $x^{q-1}\log^j x$.

If $q>k\ge0$ and $p=q-k$, then
\begin{align}
 D^{-q}P_k(x)&=
 \frac{\zeta'(1-p,x)+(H_{p-1}-H_k)\zeta(1-p,x)}{(p-1)!},
 \label{per:eq:Pprimitive}\\
 D^{-q}Q_k(x)&=
 \frac{\zeta'(1-p,x)+H_{p-1}\zeta(1-p,x)}{(p-1)!}.
 \label{per:eq:Qprimitive}
\end{align}
\end{theorem}

\begin{proof}
The normalized family gives, identically in $w$ near zero,
\begin{equation}\label{per:eq:primitivegenerator}
 D^{-q}Z_{1+w}
 =\Gamma(-w)W_{1-q+w}
 =\frac{Z_{1-q+w}}{(-w)_q}.
\end{equation}
Since
\[
 \frac1{(-w)_q}
 =-\frac1{(q-1)!w}\prod_{r=1}^{q-1}(1-w/r)^{-1},
\]
the coefficient $(-1)^nn![w^n]$ is exactly
\eqref{per:eq:primitive}. For $q\ge1$ the relevant Hurwitz jets have
order $1-q\le0$ and are locally integrable. Their means vanish, by the
integrable Hurwitz identity continued and differentiated below $1$,
so no hidden constants occur. For $q\ge2$, the shift formula isolates
the only nonanalytic endpoint term as a finite combination of
$x^{q-1}\log^j x$; it tends to zero and matches the finite endpoint
value at $1$.

For \eqref{per:eq:Pprimitive}, the Fourier coefficient after taking
$q$ primitives is $(2\pi i n)^{-p}(\gamma+L_n-H_k)$.
Differentiation of $Z_s=\Gamma(1-s)W_s$ at $s=1-p$ gives
\[
 \frac{Z'_{1-p}}{(p-1)!}
 =\partial_sW_s\big|_{s=1-p}-\psi(p)W_{1-p}.
\]
Use $\psi(p)=H_{p-1}-\gamma$. Adding the primitive of
$H_k\delta_0^{(k)}$ gives \eqref{per:eq:Qprimitive}.
\end{proof}

At $q=1$, \eqref{per:eq:primitive} becomes the familiar Stieltjes
primitive $(-1)^{n+1}\zeta^{[n+1]}(0,x)/(n+1)$. At every larger $q$
it gives all lower jets and all constants of integration explicitly.
For example,
\begin{align*}
 D^{-2}T_1(x)&=\tfrac12\zeta''(-1,x)+\zeta'(-1,x)+\zeta(-1,x),\\
 D^{-3}T_1(x)&=\tfrac14\zeta''(-2,x)+\tfrac34\zeta'(-2,x)
                                      +\tfrac78\zeta(-2,x).
\end{align*}
The family in \eqref{per:eq:Qprimitive} is the classical balanced
negative polygamma family. Equation \eqref{per:eq:Pprimitive} records
the correction required when one starts from spectral finite parts
instead of derivatives of periodic $\log\Gamma$.


\subsection{A convergent Gamma convolution as a normalization check}

One can remove the singularities by integrating. In particular, for
$g=\log\Gamma-\tfrac12\log(2\pi)$,
\begin{align}
 &\int_0^a g(x)g(a-x)\,dx
       +\int_a^1g(x)g(1+a-x)\,dx\notag\\
 &\hspace{10mm}=\zeta''(-1,a)+2\zeta'(-1,a)
                           +(2-\zeta(2))\zeta(-1,a),\quad 0<a\le1.
 \label{per:eq:gammaordinary}
\end{align}
Indeed $Dg=P_0$, so $g*g=D^{-2}(P_0*P_0)$; apply
\eqref{per:eq:00} and \eqref{per:eq:primitive}. The right side extends
continuously to the endpoints. This is a differentiated instance of
Sun's classical ordinary convolution, included to exhibit the agreement
of the singular calculus with convergent Gamma integrals.
